\setcounter{section}{1}
\setlength{\parindent}{0pt}  % niente rientro
\setlength{\parskip}{6pt}  
\section{Ruoli e Attività}

Data la composizione ridotta del team, la ripartizione dei ruoli non è stata rigida, ma flessibile ed adattabile alle necessità delle diverse fasi progettuali.
Sebbene siano stati individuati ruoli ed ambiti di responsabilità per ciascun membro, è importante sottolineare che le decisioni critiche riguardanti l'architettura del sistema e la definizione dei requisiti sono state prese congiuntamente. Inoltre, entrambi i membri hanno partecipato attivamente alle fasi di testing e revisione incrociata dei deliverable.

Di seguito sono riportati i ruoli principali assunti e le specifiche attività svolte:

\begin{table}[H]
\centering
\renewcommand{\arraystretch}{1.5} % Aumenta leggermente lo spazio tra le righe per leggibilità
\resizebox{\textwidth}{!}{
\begin{tabular}{|l|l|p{8cm}|}
\hline
\textbf{Componente del team} & \textbf{Ruolo} & \textbf{Principali attività} \\ \hline

\hline
\textbf{Enrico Sartori} & Full Stack Developer \& Software Architect & Si è occupato della definizione dell'architettura software e della scelta dello stack tecnologico. Ha gestito lo sviluppo dell'applicazione (Backend e Frontend) e la progettazione delle API. Ha inoltre curato la stesura della documentazione tecnica nei deliverable, il deployment e parte del Business Plan. \\ \hline

\textbf{Stefano Pezzetta} & System Analyst \& Frontend Developer & Si è occupato della redazione della documentazione descrittiva e dell'analisi dei requisiti. Ha gestito la modellazione del sistema (UML) e collaborato nella progettazione dello schema del Database. Ha sviluppato componenti React, contribuendo attivamente alla realizzazione e al perfezionamento dell'interfaccia utente. Ha redotto parte del Business Plan \\ \hline
\end{tabular}
}
\caption{Ruoli ricoperti e attività svolte dai membri del team}
\label{tab:ruoli_attivita}
\end{table}